\documentclass[12pt]{article}
\usepackage[T1]{fontenc}
\usepackage[utf8]{inputenc}
\usepackage{lmodern}
\usepackage{textcomp}
\usepackage{enumitem}
\usepackage{geometry}
\usepackage{graphicx}
\usepackage{amsmath, amssymb}
\geometry{margin=1in}
\usepackage{fontspec}
\setmainfont{Roboto-Regular}[Path=fonts/,Extension=.ttf]
\expandafter\newfontfamily\csname fa0000Font0\endcsname{fa0000_pos0}[
    Path=fonts/fa0000/,
    Extension=.ttf
]
\expandafter\newfontfamily\csname fa0000Font1\endcsname{fa0000_pos1}[
    Path=fonts/fa0000/,
    Extension=.ttf
]
\expandafter\newfontfamily\csname fa0000Font2\endcsname{fa0000_pos2}[
    Path=fonts/fa0000/,
    Extension=.ttf
]
\expandafter\newfontfamily\csname fa0000Font3\endcsname{fa0000_pos3}[
    Path=fonts/fa0000/,
    Extension=.ttf
]
\expandafter\newfontfamily\csname fa0000Font4\endcsname{fa0000_pos4}[
    Path=fonts/fa0000/,
    Extension=.ttf
]
\expandafter\newfontfamily\csname fa0002Font0\endcsname{fa0002_pos0}[
    Path=fonts/fa0002/,
    Extension=.ttf
]
\expandafter\newfontfamily\csname fa0002Font1\endcsname{fa0002_pos1}[
    Path=fonts/fa0002/,
    Extension=.ttf
]
\expandafter\newfontfamily\csname fa0003Font0\endcsname{fa0003_pos0}[
    Path=fonts/fa0003/,
    Extension=.ttf
]
\expandafter\newfontfamily\csname fa0003Font1\endcsname{fa0003_pos1}[
    Path=fonts/fa0003/,
    Extension=.ttf
]
\expandafter\newfontfamily\csname fa0003Font2\endcsname{fa0003_pos2}[
    Path=fonts/fa0003/,
    Extension=.ttf
]
\expandafter\newfontfamily\csname fa0004Font0\endcsname{fa0004_pos0}[
    Path=fonts/fa0004/,
    Extension=.ttf
]
\expandafter\newfontfamily\csname fa0004Font1\endcsname{fa0004_pos1}[
    Path=fonts/fa0004/,
    Extension=.ttf
]
\expandafter\newfontfamily\csname fa0004Font2\endcsname{fa0004_pos2}[
    Path=fonts/fa0004/,
    Extension=.ttf
]
\expandafter\newfontfamily\csname fa0005Font0\endcsname{fa0005_pos0}[
    Path=fonts/fa0005/,
    Extension=.ttf
]
\expandafter\newfontfamily\csname fa0005Font1\endcsname{fa0005_pos1}[
    Path=fonts/fa0005/,
    Extension=.ttf
]
\expandafter\newfontfamily\csname fa0005Font2\endcsname{fa0005_pos2}[
    Path=fonts/fa0005/,
    Extension=.ttf
]
\expandafter\newfontfamily\csname fa0005Font3\endcsname{fa0005_pos3}[
    Path=fonts/fa0005/,
    Extension=.ttf
]
\expandafter\newfontfamily\csname fa0006Font0\endcsname{fa0006_pos0}[
    Path=fonts/fa0006/,
    Extension=.ttf
]
\expandafter\newfontfamily\csname fa0006Font1\endcsname{fa0006_pos1}[
    Path=fonts/fa0006/,
    Extension=.ttf
]
\expandafter\newfontfamily\csname fa0006Font2\endcsname{fa0006_pos2}[
    Path=fonts/fa0006/,
    Extension=.ttf
]
\expandafter\newfontfamily\csname fa0006Font3\endcsname{fa0006_pos3}[
    Path=fonts/fa0006/,
    Extension=.ttf
]
\expandafter\newfontfamily\csname fa0006Font4\endcsname{fa0006_pos4}[
    Path=fonts/fa0006/,
    Extension=.ttf
]
\expandafter\newfontfamily\csname fa0006Font5\endcsname{fa0006_pos5}[
    Path=fonts/fa0006/,
    Extension=.ttf
]
\expandafter\newfontfamily\csname fa0006Font6\endcsname{fa0006_pos6}[
    Path=fonts/fa0006/,
    Extension=.ttf
]
\expandafter\newfontfamily\csname fa0006Font7\endcsname{fa0006_pos7}[
    Path=fonts/fa0006/,
    Extension=.ttf
]
\expandafter\newfontfamily\csname fa0006Font8\endcsname{fa0006_pos8}[
    Path=fonts/fa0006/,
    Extension=.ttf
]
\expandafter\newfontfamily\csname fa0007Font0\endcsname{fa0007_pos0}[
    Path=fonts/fa0007/,
    Extension=.ttf
]
\expandafter\newfontfamily\csname fa0007Font1\endcsname{fa0007_pos1}[
    Path=fonts/fa0007/,
    Extension=.ttf
]
\expandafter\newfontfamily\csname fa0007Font2\endcsname{fa0007_pos2}[
    Path=fonts/fa0007/,
    Extension=.ttf
]
\expandafter\newfontfamily\csname fa0007Font3\endcsname{fa0007_pos3}[
    Path=fonts/fa0007/,
    Extension=.ttf
]
\expandafter\newfontfamily\csname fa0007Font4\endcsname{fa0007_pos4}[
    Path=fonts/fa0007/,
    Extension=.ttf
]
\expandafter\newfontfamily\csname fa0008Font0\endcsname{fa0008_pos0}[
    Path=fonts/fa0008/,
    Extension=.ttf
]
\expandafter\newfontfamily\csname fa0008Font1\endcsname{fa0008_pos1}[
    Path=fonts/fa0008/,
    Extension=.ttf
]
\expandafter\newfontfamily\csname fa0008Font2\endcsname{fa0008_pos2}[
    Path=fonts/fa0008/,
    Extension=.ttf
]
\expandafter\newfontfamily\csname fa0009Font0\endcsname{fa0009_pos0}[
    Path=fonts/fa0009/,
    Extension=.ttf
]
\expandafter\newfontfamily\csname fa0009Font1\endcsname{fa0009_pos1}[
    Path=fonts/fa0009/,
    Extension=.ttf
]
\expandafter\newfontfamily\csname fa0009Font2\endcsname{fa0009_pos2}[
    Path=fonts/fa0009/,
    Extension=.ttf
]
\expandafter\newfontfamily\csname fa0010Font0\endcsname{fa0010_pos0}[
    Path=fonts/fa0010/,
    Extension=.ttf
]
\expandafter\newfontfamily\csname fa0010Font1\endcsname{fa0010_pos1}[
    Path=fonts/fa0010/,
    Extension=.ttf
]
\expandafter\newfontfamily\csname fa0010Font2\endcsname{fa0010_pos2}[
    Path=fonts/fa0010/,
    Extension=.ttf
]
\expandafter\newfontfamily\csname fa0010Font3\endcsname{fa0010_pos3}[
    Path=fonts/fa0010/,
    Extension=.ttf
]
\expandafter\newfontfamily\csname fa0010Font4\endcsname{fa0010_pos4}[
    Path=fonts/fa0010/,
    Extension=.ttf
]
\expandafter\newfontfamily\csname fa0010Font5\endcsname{fa0010_pos5}[
    Path=fonts/fa0010/,
    Extension=.ttf
]
\expandafter\newfontfamily\csname fa0010Font6\endcsname{fa0010_pos6}[
    Path=fonts/fa0010/,
    Extension=.ttf
]
\expandafter\newfontfamily\csname fa0010Font7\endcsname{fa0010_pos7}[
    Path=fonts/fa0010/,
    Extension=.ttf
]
\expandafter\newfontfamily\csname fa0010Font8\endcsname{fa0010_pos8}[
    Path=fonts/fa0010/,
    Extension=.ttf
]
\expandafter\newfontfamily\csname fa0010Font9\endcsname{fa0010_pos9}[
    Path=fonts/fa0010/,
    Extension=.ttf
]
\expandafter\newfontfamily\csname fa0010Font10\endcsname{fa0010_pos10}[
    Path=fonts/fa0010/,
    Extension=.ttf
]
\expandafter\newfontfamily\csname fa0010Font11\endcsname{fa0010_pos11}[
    Path=fonts/fa0010/,
    Extension=.ttf
]
\expandafter\newfontfamily\csname fa0010Font12\endcsname{fa0010_pos12}[
    Path=fonts/fa0010/,
    Extension=.ttf
]
\expandafter\newfontfamily\csname fa0011Font0\endcsname{fa0011_pos0}[
    Path=fonts/fa0011/,
    Extension=.ttf
]
\expandafter\newfontfamily\csname fa0011Font1\endcsname{fa0011_pos1}[
    Path=fonts/fa0011/,
    Extension=.ttf
]
\expandafter\newfontfamily\csname fa0011Font2\endcsname{fa0011_pos2}[
    Path=fonts/fa0011/,
    Extension=.ttf
]
\expandafter\newfontfamily\csname fa0011Font3\endcsname{fa0011_pos3}[
    Path=fonts/fa0011/,
    Extension=.ttf
]
\expandafter\newfontfamily\csname fa0011Font4\endcsname{fa0011_pos4}[
    Path=fonts/fa0011/,
    Extension=.ttf
]
\expandafter\newfontfamily\csname fa0012Font0\endcsname{fa0012_pos0}[
    Path=fonts/fa0012/,
    Extension=.ttf
]
\expandafter\newfontfamily\csname fa0012Font1\endcsname{fa0012_pos1}[
    Path=fonts/fa0012/,
    Extension=.ttf
]
\expandafter\newfontfamily\csname fa0012Font2\endcsname{fa0012_pos2}[
    Path=fonts/fa0012/,
    Extension=.ttf
]
\expandafter\newfontfamily\csname fa0012Font3\endcsname{fa0012_pos3}[
    Path=fonts/fa0012/,
    Extension=.ttf
]
\expandafter\newfontfamily\csname fa0012Font4\endcsname{fa0012_pos4}[
    Path=fonts/fa0012/,
    Extension=.ttf
]
\expandafter\newfontfamily\csname fa0012Font5\endcsname{fa0012_pos5}[
    Path=fonts/fa0012/,
    Extension=.ttf
]
\expandafter\newfontfamily\csname fa0013Font0\endcsname{fa0013_pos0}[
    Path=fonts/fa0013/,
    Extension=.ttf
]
\expandafter\newfontfamily\csname fa0013Font1\endcsname{fa0013_pos1}[
    Path=fonts/fa0013/,
    Extension=.ttf
]
\expandafter\newfontfamily\csname fa0013Font2\endcsname{fa0013_pos2}[
    Path=fonts/fa0013/,
    Extension=.ttf
]
\expandafter\newfontfamily\csname fa0013Font3\endcsname{fa0013_pos3}[
    Path=fonts/fa0013/,
    Extension=.ttf
]
\expandafter\newfontfamily\csname fa0013Font4\endcsname{fa0013_pos4}[
    Path=fonts/fa0013/,
    Extension=.ttf
]
\expandafter\newfontfamily\csname fa0013Font5\endcsname{fa0013_pos5}[
    Path=fonts/fa0013/,
    Extension=.ttf
]
\expandafter\newfontfamily\csname fa0014Font0\endcsname{fa0014_pos0}[
    Path=fonts/fa0014/,
    Extension=.ttf
]

\begin{document}
\begin{center}
Chemistry - Graduate Level Assessment 5 \
Total Marks: 40 \
Answer all questions.
\end{center}

\section*{Multiple Choice (10 marks)}
\begin{enumerate}[label=\arabic*.]
\item When a swimmer leaves {\csname fa0000Font0\endcsname\char"0077}{\csname fa0000Font1\endcsname\char"0061}{\csname fa0000Font2\endcsname\char"0072}{\csname fa0000Font3\endcsname\char"006D}{\csname fa0000Font4\endcsname\char"0020} {\char"0077}{\char"0061}{\char"0074}{\char"0065}{\char"0072} on a warm, breezy day, he experiences a cooling effect.
Why?
\begin{enumerate}[label=(\alph*)]
    \item The cold water causes a decrease in the swimmer's metabolic rate.
    \item The swimmer feels cold due to the temperature difference.
    \item Evaporation of water from the swimmer's skin cools the body down.
    \item The swimmer feels cold due to the wind chill factor.
    \item The swimmer's body is adapting to the water temperature after exiting.
\end{enumerate}
\item Determine the value of (0.0081)\^{}\ensuremath{\rule{1em}{1pt}}3/4.
\begin{enumerate}[label=(\alph*)]
    \item (100) / (27)
    \item (1000) / (81)
    \item (1000) / (9)
    \item (10) / (27)
    \item (100) / (3)
\end{enumerate}
\item What is the volume in liters of a rectangular tank which measures{\csname fa0002Font0\endcsname\char"0033}{\csname fa0002Font1\endcsname\char"002E} {\csname fa0003Font0\endcsname\char"0030}{\csname fa0003Font1\endcsname\char"0020}{\csname fa0003Font2\endcsname\char"006D} by 50 cm by 200 mm?
\begin{enumerate}[label=(\alph*)]
    \item 250 liters
    \item 450 liters
    \item 100 liters
    \item 350 liters
    \item 150 liters
\end{enumerate}
\item Using the tables of standard electrode potentials, arrange the fol-lowing substances in
decreasing order of ability as reducing agents: Al, Co, Ni, Ag, H\_2, Na .
\begin{enumerate}[label=(\alph*)]
    \item H\_2 \textgreater{} Ni \textgreater{} Co \textgreater{} Al \textgreater{} Na \textgreater{} Ag
    \item Al \textgreater{} Co \textgreater{} Ni \textgreater{} Ag \textgreater{} H\_2 \textgreater{} Na
    \item Ag \textgreater{} Na \textgreater{} Al \textgreater{} H\_2 \textgreater{} Co \textgreater{} Ni
    \item Ag \textgreater{} H\_2 \textgreater{} Ni \textgreater{} Co \textgreater{} Al \textgreater{} Na
    \item Ni \textgreater{} Co \textgreater{} Al \textgreater{} Na \textgreater{} Ag \textgreater{} H\_2
\end{enumerate}
\item A chemist wants to produce chlorine from molten NaCl. How many grams could be produced if
he uses a current of one amp for 5.00 minutes?
\begin{enumerate}[label=(\alph*)]
    \item 0.031g
    \item 5.55g
    \item 0.021g
    \item 0.511g
    \item 0.111g
\end{enumerate}
\end{enumerate}

\section*{True / False (10 marks)}
\textit{Answer True or False}
\begin{enumerate}[label=\arabic*.]
\setcounter{enumi}{5}
\item True or False: At 27°C and 60\% relative humidity, the mass of water vapor in a 400 m³ room is {\csname fa0004Font0\endcsname\char"0031}{\csname fa0004Font1\endcsname\char"0030}{\csname fa0004Font2\endcsname\char"002E} {\csname fa0005Font0\endcsname\char"0030}{\csname fa0005Font1\endcsname\char"0020}{\csname fa0005Font2\endcsname\char"006B}{\csname fa0005Font3\endcsname\char"0067}.
\item True or False: The series S {\csname fa0006Font0\endcsname\char"0063}{\csname fa0006Font1\endcsname\char"006F}{\csname fa0006Font2\endcsname\char"006E}{\csname fa0006Font3\endcsname\char"0076}{\csname fa0006Font4\endcsname\char"0065}{\csname fa0006Font5\endcsname\char"0072}{\csname fa0006Font6\endcsname\char"0067}{\csname fa0006Font7\endcsname\char"0065}{\csname fa0006Font8\endcsname\char"0073} and approaches infinity as more terms are added to the sum.
\item True or False: The pH of a 0.001 M aqueous solution of KOH is 14, which suggests it is a {\csname fa0007Font0\endcsname\char"0062}{\csname fa0007Font1\endcsname\char"0061}{\csname fa0007Font2\endcsname\char"0073}{\csname fa0007Font3\endcsname\char"0069}{\csname fa0007Font4\endcsname\char"0063} solution.
\item True or False: {\csname fa0008Font0\endcsname\char"0033}{\csname fa0008Font1\endcsname\char"0030}{\csname fa0008Font2\endcsname\char"002E} {\csname fa0009Font0\endcsname\char"0030}{\csname fa0009Font1\endcsname\char"0020}{\csname fa0009Font2\endcsname\char"0067} of Ca(OH)\_2 is required to neutralize 28.0 g of HCl, which in turn can neutralize 25.1 g of H\_3PO\_4.
\item True or False: Proteins serve as {\csname fa0010Font0\endcsname\char"006E}{\csname fa0010Font1\endcsname\char"006F}{\csname fa0010Font2\endcsname\char"006E}{\csname fa0010Font3\endcsname\char"002D}{\csname fa0010Font4\endcsname\char"0065}{\csname fa0010Font5\endcsname\char"0073}{\csname fa0010Font6\endcsname\char"0073}{\csname fa0010Font7\endcsname\char"0065}{\csname fa0010Font8\endcsname\char"006E}{\csname fa0010Font9\endcsname\char"0074}{\csname fa0010Font10\endcsname\char"0069}{\csname fa0010Font11\endcsname\char"0061}{\csname fa0010Font12\endcsname\char"006C} structural materials and biological regulators in various cellular processes.
\end{enumerate}

\section*{Long Form Response (20 marks)}
\textit{Answer each question with a clear and complete explanation.}
\begin{enumerate}[label=\arabic*.]
\setcounter{enumi}{10}
\item Given a first-order reaction with a {\csname fa0011Font0\endcsname\char"0068}{\csname fa0011Font1\endcsname\char"0061}{\csname fa0011Font2\endcsname\char"006C}{\csname fa0011Font3\endcsname\char"0066}{\csname fa0011Font4\endcsname\char"002D} {\csname fa0012Font0\endcsname\char"006C}{\csname fa0012Font1\endcsname\char"0069}{\csname fa0012Font2\endcsname\char"0066}{\csname fa0012Font3\endcsname\char"0065}{\csname fa0012Font4\endcsname\char"0020}{\csname fa0012Font5\endcsname\char"006F} {\csname fa0013Font0\endcsname\char"0066}{\csname fa0013Font1\endcsname\char"0020}{\csname fa0013Font2\endcsname\char"0036}{\csname fa0013Font3\endcsname\char"0038}{\csname fa0013Font4\endcsname\char"0020}{\csname fa0013Font5\endcsname\char"006D} {\csname fa0014Font0\endcsname\char"0069}{\char"006E}{\char"0075}{\char"0074}{\char"0065}{\char"0073}, calculate the rate constant
and discuss the relationship between half-life and rate constant for first-order
reactions.
\item Discuss the relationship between gas volume and temperature at constant pressure, using a
gas initially occupying 100 ml at 20°C to calculate its volume at 10°C, and explain the
underlying principles.
\end{enumerate}

\vfill
\noindent\textit{End of Paper}
\end{document}